\import{algebra/ring/reducts.tex}
\import{algebra/monoid/power.tex}
\import{recursion/natural.tex}

\section{Natural elements of a ring}

The natural element $n\cdot 1_R$ is the natural power of the unit in the
additive monoid of the ring.  Thus the generic monoid-power construction is
the sole recursion owner.

\begin{definition}\label{ring_natural_step}
    $\ringNaturalStepOp:\Set\to\Set\to\Set\to\Set
        =\fun{R,n,x:\Set}{
            \ringAdd[R](x,\ringOne[R])}$.
\end{definition}

\begin{definition}\label{ring_natural}
    $\ringNatural{R}{n}=
        \monoidPower{\AdditiveGroup{R}}{\ringOne[R]}{n}$.
\end{definition}

\begin{proposition}\label{ring_natural_zero}
    Let $R$ be a ring.
    Then $\ringNatural{R}{\emptyset}=\ringZero[R]$.
\end{proposition}
\begin{proof}
    Follows by
        \auto{ring_natural,monoid_power_zero,additive_group_fields}.
\end{proof}

\begin{proposition}\label{ring_natural_successor}
    Let $R$ be a ring.
    Let $n$ be a natural number.
    Then $\ringNatural{R}{\suc{n}}
        =\ringAdd[R](\ringNatural{R}{n},\ringOne[R])$.
\end{proposition}
\begin{proof}
    $\AdditiveGroup{R}$ is a monoid by
        \obv{additive_group_is_exact_commutative_group,
            commutative_group,group}.
    $\ringOne[R]\in\carrier[R]$ by \obv{ring_one_type}.
    Hence $\ringOne[R]\in\carrier[\AdditiveGroup{R}]$ by
        \auto{additive_group_fields}.
    $\ringNatural{R}{\suc{n}}=
        \mul[\AdditiveGroup{R}](
            \ringNatural{R}{n},\ringOne[R])$ by
        \auto{ring_natural,monoid_power_successor}.
    Follows by \auto{additive_group_fields}.
\end{proof}

\begin{proposition}\label{ring_natural_type}
    Let $R$ be a ring.
    Let $n$ be a natural number.
    Then $\ringNatural{R}{n}\in\carrier[R]$.
\end{proposition}
\begin{proof}
    $\AdditiveGroup{R}$ is a monoid by
        \obv{additive_group_is_exact_commutative_group,
            commutative_group,group}.
    $\ringOne[R]\in\carrier[R]$ by \obv{ring_one_type}.
    Hence $\ringOne[R]\in\carrier[\AdditiveGroup{R}]$ by
        \auto{additive_group_fields}.
    $\ringNatural{R}{n}\in\carrier[\AdditiveGroup{R}]$ by
        \auto{ring_natural,monoid_power_type}.
    Follows by \auto{additive_group_fields}.
\end{proof}

\begin{proposition}\label{ring_natural_one}
    Let $R$ be a ring.
    Then $\ringNatural{R}{\suc{\emptyset}}=\ringOne[R]$.
\end{proposition}
\begin{proof}
    $\emptyset$ is a natural number by \obv{omega_empty}.
    $\ringNatural{R}{\suc{\emptyset}}
        =\ringAdd[R](\ringNatural{R}{\emptyset},\ringOne[R])$ by
        \obv{ring_natural_successor}.
    Follows by
        \auto{ring_natural_zero,ring_one_type,ring_add_zero_left}.
\end{proof}

\begin{theorem}\label{ring_natural_addition}
    Let $R$ be a ring.
    Let $m,n$ be natural numbers.
    Then $\ringNatural{R}{m+n}=
        \ringAdd[R](\ringNatural{R}{m},\ringNatural{R}{n})$.
\end{theorem}
\begin{proof}
    $\AdditiveGroup{R}$ is a monoid by
        \obv{additive_group_is_exact_commutative_group,
            commutative_group,group}.
    $\ringOne[R]\in\carrier[R]$ by \obv{ring_one_type}.
    Hence $\ringOne[R]\in\carrier[\AdditiveGroup{R}]$ by
        \auto{additive_group_fields}.
    $\monoidPower{\AdditiveGroup{R}}{\ringOne[R]}{m+n}=
        \mul[\AdditiveGroup{R}](
            \monoidPower{\AdditiveGroup{R}}{\ringOne[R]}{m},
            \monoidPower{\AdditiveGroup{R}}{\ringOne[R]}{n})$ by
        \obv{monoid_power_addition}.
    Follows by \auto{ring_natural,additive_group_fields}.
\end{proof}

\begin{theorem}\label{ring_natural_central}
    Let $R$ be a ring.
    Let $x\in\carrier[R]$.
    Let $n$ be a natural number.
    Then $\ringMul[R](\ringNatural{R}{n},x)
        =\ringMul[R](x,\ringNatural{R}{n})$.
\end{theorem}
\begin{proof}
    Let $C=\{k\in\naturals\mid
        \ringMul[R](\ringNatural{R}{k},x)
            =\ringMul[R](x,\ringNatural{R}{k})\}$.
    $\emptyset\in C$ by
        \auto{omega_empty,ring_natural_zero,
            ring_mul_zero_left,ring_mul_zero_right}.
    Show that for every $k\in C$ we have $\suc{k}\in C$.
    \begin{subproof}
        Fix $k\in C$.
        Hence $k$ is a natural number and
            $\ringMul[R](\ringNatural{R}{k},x)
                =\ringMul[R](x,\ringNatural{R}{k})$.
        $x\in\carrier[R]$ by assumption.
        $\ringNatural{R}{k}\in\carrier[R]$ by
            \obv{ring_natural_type}.
        $\ringOne[R]\in\carrier[R]$ by \obv{ring_one_type}.
        $\ringNatural{R}{\suc{k}}=
            \ringAdd[R](\ringNatural{R}{k},\ringOne[R])$ by
            \obv{ring_natural_successor}.
        $\ringMul[R](\ringNatural{R}{\suc{k}},x)=
            \ringAdd[R](
                \ringMul[R](\ringNatural{R}{k},x),x)$ by
            \auto{ring_distributive_right,ring_mul_one_left}.
        $\ringMul[R](x,\ringNatural{R}{\suc{k}})=
            \ringAdd[R](
                \ringMul[R](x,\ringNatural{R}{k}),x)$ by
            \auto{ring_distributive_left,ring_mul_one_right}.
        $\suc{k}$ is a natural number by \obv{omega_successor}.
    \end{subproof}
    Hence $C$ is an inductive set by \obv{inductive_set}.
    Hence every natural number is an element of $C$ by
        \obv{omega_minimal}.
    Hence $n\in C$.
\end{proof}

\begin{theorem}\label{ring_natural_multiplication}
    Let $R$ be a ring.
    Let $m,n$ be natural numbers.
    Then $\ringNatural{R}{\naturalMul{m}{n}}=
        \ringMul[R](\ringNatural{R}{m},\ringNatural{R}{n})$.
\end{theorem}
\begin{proof}
    Let $C=\{k\in\naturals\mid
        \ringNatural{R}{\naturalMul{m}{k}}=
            \ringMul[R](\ringNatural{R}{m},\ringNatural{R}{k})\}$.
    $\ringNatural{R}{m}\in\carrier[R]$ by
        \obv{ring_natural_type}.
    $\emptyset\in C$ by
        \auto{omega_empty,natural_multiplication_zero_right,
            ring_natural_zero,ring_mul_zero_right}.
    Show that for every $k\in C$ we have $\suc{k}\in C$.
    \begin{subproof}
        Fix $k\in C$.
        Hence $k$ is a natural number and
            $\ringNatural{R}{\naturalMul{m}{k}}=
                \ringMul[R](\ringNatural{R}{m},
                    \ringNatural{R}{k})$.
        $\ringNatural{R}{m},\ringNatural{R}{k}\in\carrier[R]$ by
            \obv{ring_natural_type}.
        $\ringOne[R]\in\carrier[R]$ by \obv{ring_one_type}.
        $\naturalMul{m}{k}$ is a natural number by
            \obv{natural_multiplication_closed}.
        $\naturalMul{m}{\suc{k}}=\naturalMul{m}{k}+m$ by
            \obv{natural_multiplication_successor_right}.
        $\ringNatural{R}{\naturalMul{m}{\suc{k}}}=
            \ringAdd[R](
                \ringNatural{R}{\naturalMul{m}{k}},
                \ringNatural{R}{m})$ by
            \auto{ring_natural_addition}.
        $\ringNatural{R}{\suc{k}}=
            \ringAdd[R](\ringNatural{R}{k},\ringOne[R])$ by
            \obv{ring_natural_successor}.
        $\ringMul[R](\ringNatural{R}{m},
                \ringNatural{R}{\suc{k}})=
            \ringAdd[R](
                \ringMul[R](\ringNatural{R}{m},
                    \ringNatural{R}{k}),
                \ringNatural{R}{m})$ by
            \auto{ring_distributive_left,ring_mul_one_right}.
        $\suc{k}$ is a natural number by \obv{omega_successor}.
    \end{subproof}
    Hence $C$ is an inductive set by \obv{inductive_set}.
    Hence every natural number is an element of $C$ by
        \obv{omega_minimal}.
    Hence $n\in C$.
\end{proof}

\begin{theorem}\label{ring_natural_recursion_eq}
    Let $R$ be a ring.
    Let $n$ be a natural number.
    Then $\ringNatural{R}{n}
        =\Apply{\naturalRecOp}{\ringZero[R],
            \Apply{\ringNaturalStepOp}{R},n}$.
\end{theorem}
\begin{proof}
    Let $C=\{k\in\naturals\mid
        \ringNatural{R}{k}
            =\Apply{\naturalRecOp}{\ringZero[R],
                \Apply{\ringNaturalStepOp}{R},k}\}$.
    $\emptyset\in C$ by
        \auto{omega_empty,ring_natural_zero,natural_recursion_zero}.
    Show that for every $k\in C$ we have $\suc{k}\in C$.
    \begin{subproof}
        Fix $k\in C$.
        Hence $k$ is a natural number and
            $\ringNatural{R}{k}
                =\Apply{\naturalRecOp}{\ringZero[R],
                    \Apply{\ringNaturalStepOp}{R},k}$.
        $\ringNatural{R}{\suc{k}}
            =\ringAdd[R](\ringNatural{R}{k},\ringOne[R])$ by
            \obv{ring_natural_successor}.
        $\Apply{\naturalRecOp}{\ringZero[R],
                \Apply{\ringNaturalStepOp}{R},\suc{k}}
            =\ringAdd[R](
                \Apply{\naturalRecOp}{\ringZero[R],
                    \Apply{\ringNaturalStepOp}{R},k},
                \ringOne[R])$ by
            \auto{ring_natural_step,natural_recursion_successor}.
        $\suc{k}$ is a natural number by \obv{omega_successor}.
    \end{subproof}
    Hence $C$ is an inductive set by \obv{inductive_set}.
    Hence every natural number is an element of $C$ by
        \obv{omega_minimal}.
    Hence $n\in C$.
\end{proof}
